% Section 5: closure at the artificial boundaries. DtN closure and its poles; bordered closure
% (Theorem B of research/projects/bie-enclosure/theory/port-subspace.md,
% reviewed 12/12); analytic frame; real-frequency test; root search and
% acceptance (implemented in test/bie-bound/pole-free-v1/).
\section{Closure without Dirichlet poles}\label{sec:closure}

\subsection{Discrete guided modes}\label{sec:discrete-modes}
A guided mode restricted to the defect cell has data of decaying
half-strip solutions on its artificial boundaries \cite[Thm.~4.5]{Fliss2013}. The discrete counterpart is
the following.

\begin{definition}\label{def:discrete-mode}
A \emph{discrete guided mode} at $k$ is a nonzero $\xi\in\CC^{n+2K}$ with
\begin{equation}\label{eq:nep}
  \mathsf A_0(k)\xi=0,\qquad \mathsf Y_\pm(k)\xi\in\ran\mathsf Z_\pm(k).
\end{equation}
\end{definition}

For each such $\xi$ the half-strip coordinates $c_\pm$ with $\mathsf Y_\pm\xi=\mathsf Z_\pm c_\pm$
are unique, and $z_j=\mathsf Z_\pm\mathsf P_\pm^jc_\pm$ continues the field into the
half-strips. Root finding needs a scalar function of $k$ whose zeros are the $k$
at which \eqref{eq:nep} has a solution.

\subsection{The Dirichlet-to-Neumann closure and its poles}\label{sec:dtn}
The usual closure eliminates the half-strip coordinates through
Dirichlet-to-Neumann matrices. Take the normals of both artificial
boundaries pointing into the half-strips, and write the Cauchy data of the decaying
subspaces in value and derivative blocks:
\begin{equation}\label{eq:dtn}
  \mathsf C_\pm\mathsf Z_\pm=\begin{bmatrix}\mathsf D_\pm\\ \mathsf N_\pm\end{bmatrix},
  \qquad
  \Lambda_\pm=\mathsf N_\pm\mathsf D_\pm^{-1}.
\end{equation}
If the defect-cell density block $A_\eta$ is invertible, then
$H_c=\left[\begin{smallmatrix}-A_\eta^{-1}B_0\\ I\end{smallmatrix}\right]$ parametrizes
$\ker\mathsf A_0$ by the incident amplitudes $\mathsf a=(a^-_0;b^+_0)$. Let
$\mathsf T_0$ and $\mathsf N_0$ be the values and derivatives of the defect-cell field
on both artificial boundaries, $\mathsf T_0=[\,[I,I]\mathsf Y_-;[I,I]\mathsf Y_+\,]H_c$ and similarly
for $\mathsf N_0$. The closure is
\begin{equation}\label{eq:port-matrix}
  \mathsf A(k)=\mathsf N_0(k)-\diag\bigl(\Lambda_-(k),\Lambda_+(k)\bigr)\mathsf T_0(k)
  \in\CC^{2K\times2K},\qquad \det\mathsf A(k)=0 .
\end{equation}
This is the discrete analogue of the transparent condition of
\cite{Fliss2013}. It needs $\mathsf D_\pm$ invertible. The matrix $\mathsf D_\pm$
is singular exactly when a decaying half-strip solution has zero value, in the
retained orders, on the artificial boundary. This is the discrete form of a
Dirichlet exceptional value \cite[Thm.~4.1]{Fliss2013}. Its position depends
on where the artificial boundary sits in its gap strip, while the guided modes do not.

\Cref{thm:bordered} below shows that $\Lambda_\pm$ then has a genuine pole.
The determinant $\det\mathsf A$ has the factor $1/\det\mathsf D_\pm$, and zeros and
poles may cancel. Both failures of \cref{sec:task} occur in practice. In
example 1 of \cref{sec:experiments}, moving the right artificial boundary by $0.276$ puts a
pole $0.0022$ below the guided mode. There the scaled singular-value ratio
of $\mathsf A$ is $4.5\times10^{-16}$, because $\sigma_{\max}$ grows to
$2\times10^{13}$, and a ratio-based acceptance test accepts the pole as a
root. Moving it by $0.277$ puts the pole on the mode, and the search
fails.

\subsection{Bordered closure}\label{sec:bordered}
We keep the half-strip coordinates as unknowns. Let
\begin{equation}\label{eq:bordered}
  \mathsf H(k)=
  \begin{bmatrix}
    \mathsf A_0(k)&0&0\\
    \mathsf Y_-(k)&-\mathsf Z_-(k)&0\\
    \mathsf Y_+(k)&0&-\mathsf Z_+(k)
  \end{bmatrix}\in\CC^{(n+4K)\times(n+4K)},
\end{equation}
acting on $(\xi;c_-;c_+)$. The matrix is square because each decaying
subspace has dimension $K$ (\cref{lem:dimension}).

\begin{theorem}[bordered closure]\label{thm:bordered}
Let $\mathsf Z_\pm$ have full column rank $K$.
\begin{enumerate}
\item The projection $(\xi;c_-;c_+)\mapsto\xi$ is a linear isomorphism from
  $\ker\mathsf H(k)$ onto the discrete guided modes \eqref{eq:nep} at $k$. In
  particular their dimensions agree. No condition on $\mathsf D_\pm$ or on $A_\eta$
  is needed.
\item If $k$ is not a Wood point and $A_\eta$, $\mathsf D_-$, $\mathsf D_+$ are
  invertible, then
  \begin{equation}\label{eq:det-identity}
    (-1)^{K}\det\mathsf C_-\det\mathsf C_+\,\det\mathsf H
    =\det A_\eta\,\det\mathsf D_-\,\det\mathsf D_+\,\det\mathsf A .
  \end{equation}
\item Let the entries of $\mathsf H$ be holomorphic near $k_*$ in a frame of the
  kind constructed in \cref{sec:frame}. If $\det\mathsf D_\pm$ has an isolated zero
  at $k_*$, then $\Lambda_\pm$ has a pole there. If both sides of
  \eqref{eq:det-identity} are not identically zero,
  \begin{equation}\label{eq:order}
    \ord_{k_*}\det\mathsf H=\ord_{k_*}\det A_\eta+\ord_{k_*}\det\mathsf D_-
     +\ord_{k_*}\det\mathsf D_++\ord_{k_*}\det\mathsf A ,
  \end{equation}
  where $\ord$ is the order of a zero, negative for a pole.
\end{enumerate}
\end{theorem}

\begin{proof}
(1) The rows of $\mathsf H$ say $\mathsf A_0\xi=0$ and $\mathsf Y_\pm\xi=\mathsf Z_\pm c_\pm$.
Given $\xi$ with \eqref{eq:nep}, full column rank gives unique $c_\pm$.
Conversely, $\xi=0$ forces $\mathsf Z_\pm c_\pm=0$ and hence $c_\pm=0$.

(2) Multiply the two boundary block rows by $\mathsf C_-$ and $\mathsf C_+$. Then swap the
derivative rows of the left artificial boundary with the value rows of the
right one, a
permutation of sign $(-1)^{K^2}=(-1)^K$. The matrix becomes
$\left[\begin{smallmatrix}\mathsf A_0&0\\ T^c&-D\\ N^c&-N\end{smallmatrix}\right]$ with
$D=\diag(\mathsf D_-,\mathsf D_+)$, $N=\diag(\mathsf N_-,\mathsf N_+)$, and $T^c,N^c$ the
value and derivative rows of $\mathsf C_\pm\mathsf Y_\pm$. Subtract $\Lambda=ND^{-1}$
times the value rows from the derivative rows, and use $\det(-D)=\det D$ for
the even size $2K$. A block-triangular determinant is left. Eliminating the
defect-cell densities with $H_c$ turns the remaining $(n+2K)$-square block into
$\left[\begin{smallmatrix}A_\eta&*\\0&\mathsf A\end{smallmatrix}\right]$.

(3) If $\Lambda_s$ extended holomorphically to $k_*$, then
$\mathsf N_s(k_*)=\Lambda_s(k_*)\mathsf D_s(k_*)$. A nonzero
$v\in\ker\mathsf D_s(k_*)$ would give $\mathsf C_s\mathsf Z_sv=0$, contradicting full rank.
The order relation is \eqref{eq:det-identity} read as an identity of
meromorphic functions. The factors $\det\mathsf C_\pm$ do not vanish.
\end{proof}

The theorem explains both failures. The DtN determinant is the bordered
determinant divided by $\det\mathsf D_-\det\mathsf D_+$ (and by $\det A_\eta$). Its
poles are artifacts of the coordinates, while the bordered determinant has
none. At a Dirichlet exceptional value that coincides with a mode, $\det\mathsf A$
can be finite and nonzero although $\det\mathsf H$ vanishes. The scalar example
with half-strip column $(t,1)^T$ and defect-cell column $(0,1)^T$, $t=k-k_*$, has
$\Lambda=1/t$ and $\mathsf A\equiv1$, while $\det\mathsf H=t$. The same
factorization with a nonsingular Robin block in place of $\mathsf D_\pm$ gives
the Robin-to-Robin closure of \cite{FlissKlindworthSchmidt2015}. We do not
compare the two closures numerically in this paper.

\subsection{An analytic frame}\label{sec:frame}
The orthonormal QZ basis $\mathsf Z_\pm(k)$ is defined only up to a
$k$-dependent unitary factor, so $\det\mathsf H$ with the raw basis need not be
continuous in $k$. An analytic frame exists whenever $\mathsf E_{\mathrm L}-\zeta\mathsf E_{\mathrm R}$
depends holomorphically on $k$ and is invertible on a fixed contour $\Gamma$, for
example the unit circle in a gap. The Riesz projector
\begin{equation}\label{eq:riesz}
  \mathscr P(k)=\frac1{2\pi\ii}\oint_\Gamma\bigl(\zeta\mathsf E_{\mathrm R}(k)-\mathsf E_{\mathrm L}(k)\bigr)^{-1}
  \mathsf E_{\mathrm R}(k)\,\dd\zeta
\end{equation}
projects onto the decaying subspace and is holomorphic. For a fixed
reference matrix $\mathsf Z_0\in\CC^{2K\times K}$, for instance
$\mathsf Z_0=\mathsf Z_\pm(k_0)$, the matrix
\begin{equation}\label{eq:frame}
  \widehat{\mathsf Z}_\pm(k)=\mathsf Z_\pm(k)\bigl(\mathsf Z_0^*\mathsf Z_\pm(k)\bigr)^{-1}
\end{equation}
is independent of the basis chosen at each $k$. Near any $k_1$, the columns
of $F(k)=\mathscr P(k)\mathsf Z_\pm(k_1)$ form a holomorphic basis of the
decaying subspace, and \eqref{eq:frame} equals $F(k)(\mathsf Z_0^*F(k))^{-1}$,
since two bases differ by an invertible right factor that cancels. Hence
$\widehat{\mathsf Z}_\pm$ is holomorphic wherever the overlap $\mathsf Z_0^*\mathsf Z_\pm(k)$
is invertible. (The global choice $F=\mathscr P\mathsf Z_0$ does not suffice:
$\mathscr P$ is an oblique projector, and $\mathscr P(k)\mathsf Z_0$ can lose rank
while the overlap stays well conditioned.) We never evaluate \eqref{eq:riesz}:
\eqref{eq:frame} is computed from the QZ basis. A different holomorphic frame
multiplies $\det\mathsf H$ by a nonvanishing holomorphic factor, so the zeros and
their orders are unchanged.

The search works on panels of the window. On each panel, $\mathsf Z_0$ is the
QZ basis at the midpoint. The overlap is accepted if
$\sigma_{\min}(\mathsf Z_0^*\mathsf Z_\pm)\ge0.05$ and its condition number is at most
$20$ at all nodes; otherwise the panel is split. Fixed row and column scales,
computed at the midpoint, are applied to $\mathsf H$ on the whole panel.

\begin{remark}[analytic dependence of the assembly]\label{rem:holomorphy}
\Cref{thm:bordered}(3) and the frame need holomorphic matrix entries, but the
code evaluates the cell matrices only at real $k$. It suffices that the exact
discrete assembly is real-analytic on each panel: a real-analytic function on
a closed interval extends holomorphically to a complex neighborhood, with the
same zeros and orders on the interval. This holds on a closed panel, free of
Wood points, under the following conditions.
(i) The proxy matrix has constant exact rank near the panel; then the
minimum-norm fit is real-analytic. Full column rank is not available: the
proxy box of the code contains one corner point twice, so two columns
coincide for all $k$.
(ii) Every density solve in a periodicity cell and every row-joining solve used in the
assembly is nonsingular on the panel.
(iii) For each half-strip, $\mathsf E_{\mathrm L}-\zeta\mathsf E_{\mathrm R}$ is invertible on
the unit circle on the whole panel,
with $K$ multipliers inside.
(iv) The overlaps of \eqref{eq:frame} are invertible on the panel, and the
fixed scales are finite and nonzero.
The search uses floating-point samples. These carry errors from the numerical
rank decision of the least-squares solver, fast summation and rounding, and a
fixed numerical-rank threshold can jump inside a panel even when the exact
rank is constant. The tail and off-grid tests of \cref{sec:search} are
therefore empirical acceptance tests, valid only when the sample errors are
below their tolerances.
\end{remark}

\subsection{A real-frequency test}\label{sec:projection-test}
For an independent check at real $k$ we use orthogonal projections. If
$\mathsf Z_\pm^*\mathsf Z_\pm=I$ and $A_\eta$ is invertible,
\begin{equation}\label{eq:projection}
  \sigma_{\min}\begin{bmatrix}(I-\mathsf Z_-\mathsf Z_-^*)\mathsf Y_-H_c\\
                     (I-\mathsf Z_+\mathsf Z_+^*)\mathsf Y_+H_c\end{bmatrix}=0
  \quad\Longleftrightarrow\quad \ker\mathsf H(k)\ne\{0\}.
\end{equation}
Without the elimination one stacks $\mathsf A_0$ on top of the projected rows.
The left side is a nonnegative function with zeros exactly at the discrete
modes and no poles. We sample $\sigma_{\min}$ of the bordered matrix on a
dense grid to count roots independently of the determinant search.

\subsection{Root search and acceptance}\label{sec:search}
On each panel we interpolate $f(k)=\det(\text{scaled }\mathsf H(k))$, computed by
LU as log modulus and phase. The interpolation uses Chebyshev--Lobatto nodes
of degree $8,16,32,64$. A panel is resolved when the three trailing
coefficients are below $10^{-10}$ relative to the largest, and two off-grid
values agree to $10^{-9}$; otherwise it is bisected. Roots of the interpolant
come from the colleague matrix \cite{Boyd2002,ZhaoBarnett2015} and are
polished by real Newton steps.

A root at an interpolation node makes one sample nearly zero and the log
modulus very wide. Such a node is kept in the interpolant. It is proposed as
a simple zero if the two one-sided slopes agree to two percent and the
sampled value is small compared with step times slope.

A proposed root $\widehat k$ is accepted by a physical test, not by a
singular-value ratio. Take the null vector of $\mathsf H(\widehat k)$, normalize
the incident amplitudes, and reconstruct the field in the defect cell and in the
adjacent periodicity cells. Compute two jump norms at points off the quadrature
nodes: the root-mean-square jump of the traces and $x$-derivatives across
four walls, and the transmission residual on the rod boundaries. The root
must be found on two discretizations (agreement $10^{-8}$). On both, the two
jump norms must be below $0.005$, and on the finer one each must be at most
$\max(0.8\times\text{coarse},10^{-11})$. These thresholds are empirical. Poles
are excluded mainly by the structure of $\mathsf H$, not by this test: at the
former pole of the offset case the smallest singular value of $\mathsf H$
is $1.75\times10^{-3}$, against about $10^{-14}$ at roots.
